\begin{table}[H]
\centering
\caption{Welfare Losses from Demand vs.\ Supply Shocks}
\begin{threeparttable}
\small
\begin{tabular}{lcccc}
\toprule
& \multicolumn{4}{c}{CE Welfare Loss (\%)} \\
\cmidrule(lr){2-5}
Scenario & Risk-neutral & $\rho=1$ & $\rho=2$ & $\rho=5$ \\
\midrule
\multicolumn{5}{l}{\textit{Panel A: Demand shock (Great Recession analog)}} \\[3pt]
Baseline & 3.00 & 2.98 & 2.95 & 2.81 \\
No scarring ($\lambda = 0$) & 2.95 & 2.93 & 2.90 & 2.76 \\
No OLF exit ($\chi = 0$) & 2.99 & 2.97 & 2.95 & 2.84 \\
\midrule
\multicolumn{5}{l}{\textit{Panel B: Supply shock (COVID analog)}} \\[3pt]
Baseline & 0.04 & 0.0448 & 0.0473 & 0.0558 \\
\midrule
Demand/Supply ratio & 71 & 66 & 62 & 50 \\
\bottomrule
\end{tabular}
\begin{tablenotes}[flushleft]
\small
\item \textit{Notes:} Consumption-equivalent (CE) welfare losses under different risk preferences. Risk-neutral: linear utility $u(c) = c$. CRRA: $u(c) = c^{1-\rho}/(1-\rho)$; $\rho = 1$ is log utility, $\rho = 2$ is standard macro calibration, $\rho = 5$ is high risk aversion. The demand shock reduces aggregate productivity by 5\% permanently. The supply shock doubles the separation rate for 3 months. Welfare is computed over a 600-month simulation horizon with terminal-value continuation. The demand/supply ratio declines with risk aversion because CRRA compresses large consumption differences.
\end{tablenotes}
\end{threeparttable}
\label{tab:welfare}
\end{table}
